\section{Sự biến đổi quy luật giá trị thặng dư trong nền kinh tế tri thức}

\subsection{Khái quát chung: Nền kinh tế tri thức và tính nguyên vẹn của Học thuyết giá trị thặng dư}

\subsubsection{Bối cảnh chuyển đổi: Sự chuyển dịch từ nền kinh tế công nghiệp sang kinh tế tri thức và CMCN 4.0}

Nền kinh tế thế giới hiện nay đang có xu hướng dịch chuyển từ nền kinh tế công nghiệp sang nền kinh tế tri thức. Tại thời điểm C. Mác phát hiện ra học thuyết giá trị thặng dư, đối tượng kinh tế tập trung hướng đến là nguyên vật liệu, tăng trưởng kinh tế chủ yếu do nguyên liệu mang lại. Hiện nay, nền kinh tế thế giới dịch chuyển sang phát triển kinh tế tri thức, đối tượng kinh tế tri thức, chất xám, tăng trưởng kinh tế chủ yếu do tư bản vô hình mang lại \cite{tapchi_congthuong_gttd}.

Hiện nay, các nước tư bản chủ nghĩa đã và đang đạt những thành tựu nổi bật về kinh tế. Họ luôn luôn là những nước đi đầu trong việc ứng dụng những thành tựu khoa học công nghệ hiện đại đồng thời là những quốc gia đi đầu trong quá trình dịch chuyển sang nền kinh tế tri thức \cite{tapchi_congthuong_gttd}.

Bối cảnh lịch sử thay đổi thì tư duy cũng phải đổi mới, lý luận giá trị thặng dư được trình bày trong tác phẩm ``Tư bản'' của C.Mác dựa trên bối cảnh cuộc cách mạng công nghiệp lần thứ nhất, do đó có những điểm cụ thể không phù hợp với bối cảnh hiện nay, khi đang diễn ra cuộc cách mạng công nghiệp lần thứ tư và phát triển kinh tế tri thức \cite{do_the_tung_gttd}.

Trong bối cảnh kỷ nguyên mới, đặc biệt là tác động sâu rộng của Cách mạng công nghiệp lần thứ tư, kinh tế số và toàn cầu hóa, đã xuất hiện nhiều quan điểm cho rằng học thuyết giá trị thặng dư không còn phù hợp. Bước sang kỷ nguyên mới, đặc biệt là trong bối cảnh Cách mạng công nghiệp lần thứ tư, chủ nghĩa tư bản đương đại đã có nhiều điều chỉnh về hình thức sở hữu, quản lý và phân phối; khoa học, công nghệ, kinh tế số, trí tuệ nhân tạo, tự động hóa phát triển mạnh mẽ \cite{tran_hieu_gttd}.

\subsubsection{Bản chất cốt lõi không đổi: Sản xuất giá trị thặng dư là quy luật kinh tế tuyệt đối}

Mặc dù kinh tế thế giới nói chung và kinh tế các nước tư bản chủ nghĩa nói riêng đã có nhiều thay đổi, song sự thay đổi đó không làm mất đi giá trị của học thuyết giá trị thặng dư của C.Mác mà ngược lại, học thuyết giá trị thặng dư càng cần được nghiên cứu sâu hơn nữa để vận dụng vào từng điều kiện thực tiễn của đất nước. Trong xu hướng phát triển kinh tế tri thức hiện nay, học thuyết giá trị thặng dư của C.Mác vẫn giữ nguyên giá trị \cite{tapchi_congthuong_gttd}.

Những nguyên lý cơ bản về giá trị thặng dư vẫn giữ nguyên tính khoa học, chỉ có những biểu hiện cụ thể biến đổi \cite{do_the_tung_gttd}. Như vậy, chuyển đổi nền kinh tế từ công nghiệp sang nền kinh tế tri thức không làm mất đi giá trị của học thuyết giá trị thặng dư của C.Mác mà thậm chí, mục tiêu của các nhà tư bản khi chuyển đối tượng kinh tế sang đầu tư các loại hàng hóa chiếm hàm lượng tri thức cao cũng đều phục vụ nhu cầu tìm kiếm giá trị thặng dư và lợi nhuận siêu ngạch \cite{tapchi_congthuong_gttd}.

Học thuyết giá trị thặng dư của C. Mác là hạt nhân lý luận của kinh tế chính trị Mác -Lênin, giữ vai trò then chốt trong việc vạch rõ bản chất, quy luật vận động và xu hướng phát triển của chủ nghĩa tư bản. Học thuyết giá trị thặng dư đã chỉ ra một cách khoa học rằng, bóc lột trong chủ nghĩa tư bản không phải là hiện tượng ngẫu nhiên hay đạo đức cá nhân, mà là quan hệ kinh tế khách quan, gắn chặt với chế độ chiếm hữu tư nhân tư bản chủ nghĩa về tư liệu sản xuất \cite{tran_hieu_gttd}.

Giá trị thặng dư là quy luật kinh tế tuyệt đối của phương thức sản xuất tư bản chủ nghĩa, không có sản xuất giá trị thặng dư thì không có chủ nghĩa tư bản, giá trị thặng dư là nguồn gốc của mâu thuẫn cơ bản, nội tại của xã hội tư bản. Chừng nào còn chế độ chiếm hữu tư bản tư nhân về tư liệu sản xuất, còn hàng hóa sức lao động, chừng nào mà người lao động còn phải thêm vào thời gian lao động cần thiết để nuôi sống mình một số thời gian lao động dôi ra để sản xuất những tư liệu sinh hoạt cho người chiếm hữu tư liệu sản xuất, chừng đó, học thuyết giá trị thặng dư của C. Mác vẫn còn nguyên giá trị \cite{tran_hieu_gttd}.

\subsubsection{Khẳng định nguyên lý: Máy móc không tự tạo ra GTTD, chỉ có lao động sống mới tạo ra}

Từ đó, xuất hiện những luận điểm cho rằng học thuyết giá trị thặng dư đã lỗi thời, rằng ngày nay ``máy móc tạo ra giá trị'', hoặc ``chủ nghĩa tư bản chỉ còn bóc lột người máy'' \cite{tran_hieu_gttd}. Có những máy móc tự động không có người điều khiển, dường như lao động sống bằng không. Bởi vậy, có người cho rằng máy móc tạo ra của cải không cần sức lao động, nên máy móc tạo ra giá trị thặng dư \cite{do_the_tung_gttd}.

Nhưng vẫn phải có người chế tạo và lập trình cho máy móc tự động. C.Mác đã giải thích: \textbf{Thiên nhiên không chế tạo ra máy móc, tất cả những máy móc đều là sản phẩm lao động của con người, đều là những cơ quan của bộ óc con người do bàn tay con người tạo ra, đều là sức mạnh đã vật hóa của tri thức} \cite{do_the_tung_gttd, marx_toantp_t46p2}.

Những lập luận đó tuy có sức hấp dẫn bề ngoài, nhưng không bác bỏ được bản chất khoa học của học thuyết. Thực tế cho thấy, máy móc, trí tuệ nhân tạo, thuật toán hay dữ liệu lớn không tự sản sinh ra giá trị mới. Chúng đều là sản phẩm của lao động quá khứ, thuộc về tư liệu sản xuất, và chỉ chuyển giá trị của mình vào sản phẩm, chứ không tạo ra giá trị thặng dư \cite{tran_hieu_gttd}.

Giá trị thặng dư vẫn là một bộ phận giá trị mới dư ra ngoài giá trị sức lao động do lao động sống tạo ra, còn lao động quá khứ kết tinh trong tư liệu sản xuất chỉ chuyển giá trị sang sản phẩm theo mức độ tiêu dùng, không tạo ra giá trị thặng dư \cite{do_the_tung_gttd}.

Như vậy, điểm mấu chốt của học thuyết giá trị thặng dư là: chỉ có lao động sống mới tạo ra giá trị của hàng hóa, tạo ra giá trị thặng dư. Nguồn gốc tạo ra giá trị thặng dư là sức lao động của công nhân làm thuê, chỉ có lao động sống (sức lao động đang hoạt động) mới tạo ra giá trị \cite{tran_hieu_gttd}.

Dù những hình thức bóc lột có biến tướng, tinh vi đến mức nào chăng nữa thì bản chất bóc lột của tư bản đối với lao động vẫn là bóc lột giá trị thặng dư -- tức là bóc lột lao động sống của người lao động chứ không thể bóc lột lao động ``chết'' của máy móc được. Do vậy, nói một cách khác, học thuyết giá trị thặng dư vẫn còn nguyên giá trị, chủ nghĩa tư bản vẫn giữ nguyên bản chất bóc lột của nó \cite{tran_hieu_gttd}.

\subsection{Sự biến đổi về nguồn gốc và động lực tạo ra giá trị thặng dư}

\subsubsection{Sự dịch chuyển hình thái lao động: Chuyển từ lao động cơ bắp sang bóc lột hàm lượng tri thức, chất xám}

Cơ cấu lao động xã hội ở các nước tư bản phát triển hiện nay có sự biến đổi lớn. Hàm lượng chất xám (sự đầu tư trí tuệ của công nhân, trí thức, kỹ thuật lập trình, nhà khoa học phát minh, sáng chế, nhà quản trị và công nghệ hiện đại), nên lao động phức tạp, lao động trí tuệ tăng lên và thay thế lao động giản đơn, lao động cơ bắp, dẫn đến lao động trí tuệ, lao động có trình độ kỹ thuật cao ngày càng có vai trò quyết định trong việc sản xuất ra giá trị thặng dư \cite{tran_hieu_gttd}.

Việc tự động hóa, số hóa, nền tảng hóa sản xuất không làm mất đi vai trò của lao động sống, mà thay đổi hình thức biểu hiện của nó: từ lao động cơ bắp sang lao động trí tuệ, lao động sáng tạo, lao động dữ liệu, lao động nền tảng số \cite{tran_hieu_gttd}.

Đáng lưu ý là trong nền kinh tế tri thức, yếu tố quan trọng nhất mang lại nguồn lượng giá trị thặng dư lớn cho nền kinh tế tư bản không phải là tư bản hữu hình, mà là tư bản vô hình tồn tại trong hàng hóa sức lao động. Những người lao động ở đây là người lao động có tri thức và trình độ, là những người chủ sở hữu kinh tế tiềm năng -- tiềm năng về trình độ, về tri thức khoa học \cite{tapchi_congthuong_gttd}.

Sự thay đổi này thực chất là thay đổi hình thức bóc lột giá trị thặng dư của nhà tư bản đối với người lao động. Chuyển bóc lột lao động từ hao phí lao động chân tay sang bóc lột hàm lượng tri thức, chất xám, gia tăng khối lượng giá trị thặng dư từ lao động của người công nhân nhưng lại có được sự quan tâm và nhiệt huyết của họ, đồng thời xóa mờ mâu thuẫn giữa hai giai cấp đối kháng này \cite{tapchi_congthuong_gttd}.

Bất kỳ lao động nào -- dù là lập trình, thiết kế thuật toán, quản trị dữ liệu, sáng tạo nội dung số -- nếu không được trả công tương xứng với giá trị mà nó tạo ra, thì đều là đối tượng bị bóc lột giá trị thặng dư. Thu nhập chủ yếu của nhà tư bản trong điều kiện kinh tế tri thức cũng không phải công quản lý mà từ phần lao động thặng dư của người lao động làm thuê, chủ yếu là lao động trí tuệ, nhà tư bản chiếm lấy \cite{tran_hieu_gttd}.

\subsubsection{Khoa học công nghệ trở thành lực lượng sản xuất trực tiếp: Nguồn gốc chủ yếu tạo ra của cải}

Trong cuộc cách mạng công nghiệp lần thứ nhất, lao động sống trực tiếp của quá trình lao động là nguồn gốc chủ yếu tạo ra giá trị thặng dư. Lao động khoa học nằm ngoài quá trình sản xuất và khoảng cách giữa khoa học, công nghệ và sản xuất xa nhau. Còn trong kinh tế tri thức và trong cách mạng công nghiệp lần thứ tư thì khoa học trở thành lực lượng sản xuất trực tiếp, tri thức và lao động khoa học -- công nghệ trở thành nguồn gốc chủ yếu của của cải và của giá trị thặng dư \cite{do_the_tung_gttd}.

C.Mác đã dự báo: \textbf{Theo đà phát triển của đại công nghiệp, việc tạo ra của cải thực sự trở nên ít phụ thuộc vào thời gian lao động và số lượng lao động đã chi phí mà chúng phụ thuộc vào trình độ chung của khoa học và vào tiến bộ kỹ thuật, hay là phụ thuộc vào việc ứng dụng khoa học ấy vào sản xuất} \cite{do_the_tung_gttd, marx_toantp_t46p2}.

Lao động biểu hiện ra không phải chủ yếu là lao động được nhập vào quá trình sản xuất mà chủ yếu là một loại lao động, trong đó con người là người kiểm soát và điều tiết bản thân quá trình sản xuất. Hệ thống máy móc tự động sẽ từng bước thay thế hầu hết lao động trực tiếp. Bởi vậy thay vì làm tác nhân chủ yếu của quá trình sản xuất, công nhân lại đứng bên cạnh quá trình ấy \cite{do_the_tung_gttd, marx_toantp_t46p2}.

Khi ấy, tri thức trở thành lực lượng sản xuất trực tiếp, phát minh trở thành một nghề đặc biệt và đối với nghề này thì việc vận dụng khoa học vào nền sản xuất trực tiếp tự nó trở thành một trong những yếu tố có tính chất quyết định và kích thích. Quá trình sản xuất từ chỗ là một quá trình lao động giản đơn thành một quá trình khoa học \cite{do_the_tung_gttd, marx_toantp_t46p2}.

Các yếu tố vật chất thì cùng với sự tiến bộ không ngừng của khoa học -- công nghệ cũng được thay thế bằng những cái mới có hiệu quả hơn khiến cho một khối lượng lao động ít hơn (nhưng chất lượng cao hơn) cũng có thể vận dụng một khối lượng máy móc và nguyên liệu lớn hơn \cite{do_the_tung_gttd, marx_toantp_t46p2}.

\subsection{Sự biến đổi về phương thức và hình thức bóc lột (Ngày càng tinh vi hơn)}

\subsubsection{Sự che giấu qua các ``van điều áp'' xã hội}

Sự xuất hiện chế độ sở hữu hỗn hợp với sự hiện diện của các công ty cổ phần, trong đó đại bộ phận là sở hữu tư nhân tư bản với một bộ phận nhỏ cổ phần của người lao động đã làm giảm đi một phần nào tính gay gắt của mâu thuẫn giữa tính chất xã hội hóa cao của lực lượng sản xuất xã hội với chế độ sở hữu tư nhân tư bản chủ nghĩa về tư liệu sản xuất. Việc cho người công nhân được mua cổ phiếu, tham dự hội nghị cổ đông, việc giảm thiểu thời gian lao động trong tuần,... dường như là chiếc van điều áp, giảm thiểu mâu thuẫn giữa tư bản và lao động \cite{tran_hieu_gttd}.

Ngày nay, sự điều tiết phân phối giá trị thặng dư của các nhà tư bản qua thuế, quỹ phúc lợi xã hội, bảo hiểm xã hội, trợ cấp thất nghiệp,... cũng tạo nên một số thu nhập nào đó cho người lao động \cite{tran_hieu_gttd}.

Giai cấp tư sản dưới danh nghĩa bán 1\% cổ phần cho người lao động, họ đã thực hiện được phương châm ``bỏ con săn sắt, bắt con cá rô'', họ chia tỉ lệ 1\% giá trị thặng dư cho người lao động, nhận về 99\% lao động không công do công nhân làm ra đồng thời kèm theo đó là sự quan tâm, tinh thần thái độ lao động hết mình của người lao động. Người công nhân lại làm việc trên tinh thần ``làm cho mình'' với nhiệt huyết hăng say, năng suất và hiệu quả tăng lên \cite{tapchi_congthuong_gttd}.

Hiện nay, trong phạm vi quốc gia, chủ nghĩa tư bản hiện đại cố gắng xây dựng một hệ thống pháp luật đa dạng, phổ cập trên các lĩnh vực của đời sống xã hội. Vì thế, việc nhà nước tư sản ở các nước công nghiệp phát triển chiếm hữu và phân phối từ 30\% -- 60\% thu nhập quốc dân và sử dụng một phần từ siêu lợi nhuận thu được để trả công cho người lao động dễ tạo ra trong người lao động một ``ảo giác'' về tình trạng không bị bóc lột \cite{tran_hieu_gttd}.

\subsubsection{Bóc lột trong không gian số và kinh tế nền tảng (Gig Economy)}

Đặc điểm đặc thù của bóc lột tư bản chủ nghĩa là quan hệ bóc lột bị che giấu sau quan hệ giữa vật với vật, giữa tiền với tiền, giữa tư bản với tư bản. Tỷ lệ bóc lột không còn hiện ra như tỷ lệ giữa người với người, mà biến thành tỷ lệ giữa giá trị thặng dư và tư bản khả biến. Chính sự vật hóa này làm cho bóc lột trong chủ nghĩa tư bản trở nên tinh vi, khó nhận diện nhưng không giới hạn \cite{tran_hieu_gttd}.

Trong kỷ nguyên số, bóc lột không giảm đi, mà trở nên tinh vi hơn, ẩn dưới các hợp đồng linh hoạt, làm việc tự do, kinh tế nền tảng, và sự phụ thuộc vào các tập đoàn công nghệ xuyên quốc gia. Dù những hình thức bóc lột có biến tướng, tinh vi đến mức nào chăng nữa thì bản chất bóc lột của tư bản đối với lao động vẫn là bóc lột giá trị thặng dư \cite{tran_hieu_gttd}.

\subsubsection{Công nghệ san bằng lợi thế tự nhiên}

Sự tiến bộ của nền sản xuất xã hội có tác dụng san bằng đối với độ mầu mỡ và vị trí của các khoảnh đất với tư cách là cơ sở của địa tô chênh lệch. Những điều kiện thuận lợi ấy do lực lượng tự nhiên mang lại không phải là nguồn gốc sinh ra lợi nhuận siêu ngạch, mà chỉ là cơ sở tự nhiên của lợi nhuận siêu ngạch \cite{do_the_tung_gttd, marx_toantp_t25p2}.

Nhưng sự phát triển của khoa học và công nghệ đã có tác dụng san bằng sự khác nhau về vị trí và độ phì nhiêu của các khoảnh đất. Mức độ phì nhiêu thực tế tùy thuộc vào sự phát triển và ứng dụng khóa học và công nghệ vào sản xuất \cite{do_the_tung_gttd, marx_toantp_t25p2}. Một thí dụ điển hình là Israel, một nước có tới 70\% lãnh thổ là sa mạc, khô cằn, thiếu nước ngọt, nhưng nhờ coi trọng nghiên cứu và ứng dụng các thành tựu khoa học và công nghệ đã có một nền nông nghiệp thuộc loại hàng đầu thế giới \cite{do_the_tung_gttd}.

\subsection{Sự biến đổi về vị trí của người lao động và thước đo sự giàu có}

\subsubsection{Vai trò mới của người lao động}

Chủ nghĩa tư bản hiện đã có nhiều thay đổi đáng kể về quan hệ sở hữu, về vai trò của người lao động trong doanh nghiệp (người lao động trở thành những cổ đông đồng sở hữu về tư liệu sản xuất). Đặc biệt, trong nền kinh tế tri thức, vai trò của người lao động ngày càng được coi trọng, những sáng kiến kinh nghiệm, những phát minh khoa học của người lao động được chủ tư bản đánh giá cao \cite{tapchi_congthuong_gttd}.

Lao động biểu hiện ra không phải chủ yếu là lao động được nhập vào quá trình sản xuất mà chủ yếu là một loại lao động, trong đó con người là người kiểm soát và điều tiết bản thân quá trình sản xuất. Bởi vậy thay vì làm tác nhân chủ yếu của quá trình sản xuất, công nhân lại đứng bên cạnh quá trình ấy. Lao động trực tiếp về lượng sẽ quy vào một phần nhỏ hơn, còn về chất được chuyển hóa thành một yếu tố cần thiết \cite{do_the_tung_gttd, marx_toantp_t46p2}.

Dù máy móc tự động hóa đã thay thế lao động trực tiếp và một bộ phận lao động trí óc của con người, nhưng không thay thế được địa vị lao động của con người, càng không thể thay đổi địa vị chủ thể của con người trong quá trình sản xuất \cite{tran_hieu_gttd}.

\subsubsection{Thước đo giá trị dịch chuyển (Dự báo của C.Mác)}

C.Mác cũng dự báo: \textbf{Một khi lao động dưới hình thái trực tiếp không còn là nguồn gốc của của cải nữa, thì nền sản xuất dựa trên giá trị trao đổi sẽ bị sụp đổ và sự rút ngắn thời gian lao động cần thiết nhằm tăng thời gian lao động thặng dư sẽ thay bằng tăng thời gian nhàn rỗi cho xã hội nói chung và cho từng thành viên của xã hội, nghĩa là tạo ra khả năng rộng rãi để phát triển một cách hoàn toàn đầy đủ lực lượng sản xuất của từng người} \cite{do_the_tung_gttd, marx_toantp_t46p2_b}.

Nếu sức sản xuất của lao động tăng lên nhờ ứng dụng các thành tựu khoa học -- công nghệ hiện đại vẫn nhằm vào mục tiêu rút ngắn thời gian lao động cần thiết để tăng lao động thặng dư thì tình trạng thất nghiệp và sản xuất thừa tương đối sẽ ngày càng trầm trọng. Khi nào công nhân bắt đầu thực hiện việc đó, thì khi ấy, một mặt thước đo thời gian lao động cần thiết sẽ là những nhu cầu của cá nhân xã hội \cite{do_the_tung_gttd, marx_toantp_t46p2_b}.

Khi ấy thước đo sự giàu có tuyệt nhiên sẽ không còn là thời gian lao động nữa, mà là thời gian nhàn rỗi. Tiết kiệm thời gian lao động đồng nghĩa với tăng thời gian nhàn rỗi, nghĩa là thời gian cho sự phát triển đầy đủ của cá nhân, đến lượt nó, bản thân sự phát triển ấy, với tính cách là sức sản xuất vô cùng to lớn, tác động trở lại sức sản xuất của lao động \cite{do_the_tung_gttd, marx_toantp_t46p2_b}.
